\chapter{2005年上海卷（秋文）}
\section{填空题}
\begin{problem}
函数 \( f(x)=\log_{4}(x+1) \) 的反函数 \( f^{-1}(x)= \)\fillinblank{}.
\end{problem}

\begin{problem}
方程 \( 4^{x} + 2^{x} - 2 = 0 \) 的解是\fillinblank{}.
\end{problem}

\begin{problem}
若 \(x, y\) 满足条件 \(\left\{\begin{aligned} &x + y \leqslant 3, \\ &y \leqslant 2x, \end{aligned}\right.\) 则 \(z = 3x + 4y\) 的最大值是\fillinblank{}.
\end{problem}

\begin{problem}
直角坐标平面 \(xOy\) 中，若定点 \(A(1,2)\) 与动点 \(P(x,y)\) 满足 \(\overrightarrow{OP} \cdot \overrightarrow{OA} = 4\)，则点 \(P\) 的轨迹方程是\fillinblank{}.
\end{problem}

\begin{problem}
函数 \( y = \cos 2x + \sin x \cos x \) 的最小正周期 \( T = \)\fillinblank{}.
\end{problem}

\begin{problem}
若 \( \cos\alpha=\frac{1}{7} \)，\(\alpha\in\left(0,\frac{\pi}{2}\right)\)，则 \( \cos\left(\alpha+\frac{\pi}{3}\right)= \)\fillinblank{}.
\end{problem}

\begin{problem}
若椭圆长轴长与短轴长之比为 \(2\)，它的一个焦点是 \( (2\sqrt{15},0) \)，则椭圆的标准方程是\fillinblank{}.

\end{problem}

\begin{problem}
某班有 \(50\) 名学生，其中 \(15\) 人选修 A 课程，另外 \(35\) 人选修 B 课程. 从班级中任选两名学生，他们是选修不同课程的学生的概率是\fillinblank{}. （结果用分数表示）
\end{problem}

\begin{problem}
直线 \( y = \frac{1}{2} x \) 关于直线 \( x = 1 \) 对称的直线方程是\fillinblank{}.
\end{problem}

\begin{problem}
在 \(\triangle ABC\) 中，若 \(A = 120^{\circ}\)，\(AB = 5\)，\(BC = 7\)，则 \(AC = \)\fillinblank{}.
\end{problem}

\begin{problem}
函数 \( f(x)=\sin x+2|\sin x| \)，\( x\in[0,2\pi] \) 的图象与直线 \( y=k \) 有且仅有两个不同的交点，则 \( k \) 的取值范围是\fillinblank{}.
\end{problem}

\begin{problem}
有两个相同的直三棱柱，高为 \(\frac{2}{a}\)，底面三角形的三边长分别为 \(3a, 4a, 5a\) \((a > 0)\). 用它们拼成一个三棱柱或四棱柱，在所有可能的情形中，全面积最小的是一个四棱柱，则 \(a\) 的取值范围是\fillinblank{}.

\bitmapfigure[max width=0.92\linewidth,max totalheight=4.2cm]{img/2005/shanghai_liberal/q12_fig1.png}
\end{problem}

\section{单选题}
\begin{problem}
若函数 \( f(x) = \frac{1}{2^x + 1} \)，则该函数在 \((- \infty, + \infty)\) 上是（\quad）

\choices
    {单调递减无最小值}
    {单调递减有最小值}
    {单调递增无最大值}
    {单调递增有最大值}
\end{problem}

\begin{problem}
已知集合 \(M = \{x \mid |x - 1| \leqslant 2, x \in \R\}\)，\(P = \left\{x \mid \frac{5}{x + 1} \geqslant 1, x \in \Z\right\}\)，则 \(M \cap P\) 等于（\quad）

\choices
    {\(\{x\mid 0 < x\leqslant 3,x\in \Z\}\)}
    {\(\{x\mid 0\leqslant x\leqslant 3,x\in \Z\}\)}
    {\(\{x\mid -1\leqslant x\leqslant 0,x\in \Z\}\)}
    {\(\{x\mid -1\leqslant x < 0,x\in \Z\}\)}
\end{problem}

\begin{problem}
条件甲：“\(a > 1\)”是条件乙：“\(a > \sqrt{a}\)”的（\quad）

\choices
    {既不充分也不必要条件}
    {充要条件}
    {充分不必要条件}
    {必要不充分条件}
\end{problem}

\begin{problem}
用 \(n\) 个不同的实数 \(a_1, a_2, \dots, a_n\) 可得到 \(n!\) 个不同的排列，每个排列为一行写成一个 \(n!\) 行的数阵. 对第 \(i\) 行 \(a_{i1}, a_{i2}, \dots, a_{in}\)，记 \(b_i = -a_{i1} + 2a_{i2} - 3a_{i3} + \dots + (-1)^n n a_{in}\)，\(i = 1, 2, 3, \dots, n!\). 例如：用 \(1, 2, 3\) 可得数阵如图，由于数阵中每一列各数之和都是 \(12\)，所以，\(b_1 + b_2 + \cdots + b_6 = -12 + 2 \times 12 - 3 \times 12 = -24\)，那么，在用 \(1, 2, 3, 4, 5\) 形成的数阵中，
\(b_1 + b_2 + \cdots + b_{120} = \)\fillinblank{}.

\begin{center}
\begin{tabular}{ccc}
1 & 2 & 3 \\
1 & 3 & 2 \\
2 & 1 & 3 \\
2 & 3 & 1 \\
3 & 1 & 2 \\
3 & 2 & 1
\end{tabular}
\end{center}

\choices
    {\(-3600\)}
    {1800}
    {\(-1080\)}
    {\(-720\)}
\end{problem}

\section{解答题}
\begin{problem}
已知长方体 \(ABCD-A_{1}B_{1}C_{1}D_{1}\) 中，\(M\)，\(N\) 分别是 \(BB_{1}\) 和 \(BC\) 的中点，\(AB = 4\)，\(AD = 2\)，\(B_{1}D\) 与平面 \(ABCD\) 所成角的大小为 \(60^{\circ}\)，求异面直线 \(B_{1}D\) 与 \(MN\) 所成角的大小. （结果用反三角函数值表示）
\bitmapfigure[max width=0.42\linewidth,max totalheight=4.8cm]{img/2005/shanghai_liberal/q17_fig1.png}
\end{problem}

\begin{problem}
在复数范围内解方程 \(|z|^2 + (z + \overline{z})\i = \frac{3 - \i}{2 + \i}\). （\(\i\) 为虚数单位）
\end{problem}

\begin{problem}
已知函数 \( f(x) = kx + b \) 的图象与 \(x, y\) 轴分别相交于点 \(A, B\)，\(\overrightarrow{AB} = 2\overrightarrow{i} + 2\overrightarrow{j}\)（\(\overrightarrow{i}, \overrightarrow{j}\) 分别是与 \(x, y\) 轴正半轴同方向的单位向量），函数 \( g(x) = x^2 - x - 6 \).

\begin{enumerate}
\item 求 \(k, b\) 的值；
\item 当 \(x\) 满足 \(f(x) > g(x)\) 时，求函数 \(\frac{g(x) + 1}{f(x)}\) 的最小值.
\end{enumerate}
\end{problem}

\begin{problem}
假设某市 2004 年新建住房面积 \(400\) 万平方米，其中有 \(250\) 万平方米是中低价房. 预计在今后的若干年内，该市每年新建住房面积平均比上一年增长 \(8\%\). 另外，每年新建住房中，中低价房的面积均比上一年增加 \(50\) 万平方米. 那么，到哪一年底，

\begin{enumerate}
\item 该市历年所建中低价房的累计面积（以 2004 年为累计的第一年）将首次不少于 \(4750\) 万平方米？
\item 当年建造的中低价房的面积占该年建造住房面积的比例首次大于 \(85\%\)？
\end{enumerate}
\end{problem}

\begin{problem}
已知抛物线 \( y^{2} = 2px \) \((p > 0)\) 的焦点为 \(F\)，\(A\) 是抛物线上横坐标为 \(4\)，且位于 \(x\) 轴上方的点，\(A\) 到抛物线准线的距离等于 \(5\). 过 \(A\) 作 \(AB\) 垂直于 \(y\) 轴，垂足为 \(B\)，\(OB\) 的中点为 \(M\).

\begin{enumerate}
\item 求抛物线方程；
\item 过 \(M\) 作 \(MN \perp FA\)，垂足为 \(N\)，求点 \(N\) 的坐标；
\item 以 \(M\) 为圆心，\(MB\) 为半径作圆 \(M\)，当 \(K(m,0)\) 是 \(x\) 轴上一动点时，讨论直线 \(AK\) 与圆 \(M\) 的位置关系.
\end{enumerate}

\bitmapfigure[max width=0.42\linewidth,max totalheight=5.2cm]{img/2005/shanghai_liberal/q21_fig1.png}

\end{problem}

\begin{problem}
对定义域是 \(D_f\)，\(D_g\) 的函数 \(y = f(x)\)，\(y = g(x)\)，规定：函数
\(h(x) = \begin{cases}
f(x)g(x), & \text{当 } x \in D_f \text{ 且 } x \in D_g, \\
f(x), & \text{当 } x \in D_f \text{ 且 } x \notin D_g, \\
g(x), & \text{当 } x \notin D_f \text{ 且 } x \in D_g.
\end{cases}\)

\begin{enumerate}
\item 若函数 \( f(x) = -2x + 3\)，\(x \geqslant 1\)，\(g(x) = x - 2\)，\(x \in \R\)，写出函数 \(h(x)\) 的解析式；
\item 求问题（1）中函数 \(h(x)\) 的最大值；
\item 若 \(g(x) = f(x + \alpha)\)，其中 \(\alpha\) 是常数，且 \(\alpha \in [0,\pi]\)，请设计一个定义域为 \(\R\) 的函数 \(y = f(x)\)，及一个 \(\alpha\) 的值，使得 \(h(x) = \cos 4x\)，并予以证明.
\end{enumerate}
\end{problem}
